\section{Surface waves: thermoelastic Rayleigh waves}\label{sec:surface}
\subsection{Formulation}
Consider the half-space $y>0$ in plane strain, with waves propagating along $x$, $\propto e^{\ii kx}$ and decaying with depth. With potentials $\bu=\nabla\phi+\nabla\times(\psi\bm e_z)$, each coupled dilatational--thermal mode has $\phi_j,\theta_j\propto e^{-m_jy}$ with
\begin{equation}
m_j^2=k^2+\kappa_j,\qquad G\kappa^2-\bigl(Gs^2+s+\varepsilon s\bigr)\kappa+s^3=0\quad(\rho=c=c_1=T_0=1),
\end{equation}
which is \eqref{eq:disp} with $K=-\kappa$, and the shear potential decays with $m_s^2=k^2+s^2/c_2^2$. The ratio of temperature to potential amplitude in each coupled mode is $\theta_j=(\kappa_j-s^2)\phi_j/\beta$. The boundary conditions on $y=0$ are traction-free, $\sigma_{yy}=\sigma_{xy}=0$, and either an insulated ($\partial_y\theta=0$, which is $q_y=0$ for any $G$) or an isothermal ($\theta=0$) thermal condition. With
\begin{align}
\sigma_{yy}&=\sum_{j=1}^2\bigl[\lambda\kappa_j+2\mu m_j^2-\beta\theta_j/\phi_j\bigr]\phi_j+2\mu\,\ii k\,m_s\,\psi,\\
\sigma_{xy}&=\mu\Bigl[-2\ii k\sum_{j=1}^2m_j\phi_j+(m_s^2+k^2)\psi\Bigr],
\end{align}
the three conditions form a homogeneous $3\times3$ linear system for $(\phi_1,\phi_2,\psi)$ whose determinant $\mathcal D(s,k)$ must vanish. The root $s(k)$ with $\Real m_j>0$ gives the complex Rayleigh frequency $\omega=\ii s$ for real $k$ (temporal damping); the same determinant at real $\omega$ and complex $k$ gives spatial damping. Differences between formulations in the literature (which wave attenuation measure is plotted, which branch of the square roots is used, how the thermal boundary condition is chosen) account for many apparent disagreements among published Rayleigh dispersion curves \citep{Biswas2020a,Tung2019,Gupta2024,Aldandani2026,Abouelregal2011,Biswas2020c,Chirita2012,Singh2013,Singh2015}.

\paragraph{Check of the limit.} For $\varepsilon\to0$ the determinant factorises and the root reduces to the classical Rayleigh equation $(2-c^2/c_2^2)^2=4\sqrt{1-c^2/c_1^2}\sqrt{1-c^2/c_2^2}$ \citep{Rayleigh1885}. For $\nu=0.3$ (so $c_2^2/c_1^2=(1-2\nu)/(2-2\nu)=2/7$) the root is $c_R=0.4957\,c_1=0.9274\,c_2$; the numerical determinant root at $\varepsilon=10^{-4}$ reproduces this value to the digits displayed, and the thermoelastic root approaches it at large $k$ where the response becomes isothermal.

\begin{figure}[t!]\centering\small
% fig5_rayleigh.tex -- Rayleigh wave geometry on half-space (s10_surface)
\begin{tikzpicture}[x=1.0cm,y=1.0cm,>=Stealth]
  \tikzset{
    node/.style={draw,rounded corners=1pt,fill=blue!7,align=center,font=\footnotesize,inner sep=2pt,text width=2.4cm},
    arr/.style={->,thick,shorten >=1pt,shorten <=1pt},
    lbl/.style={font=\footnotesize,text=blue!60!black},
    medium/.style={fill=blue!10},
  }
  % Half-space region
  \fill[medium] (-3.5,-3.5) -- (3.5,-3.5) -- (3.5,0) -- (-3.5,0) -- cycle;
  % Surface line
  \draw[thick] (-3.5,0) -- (3.5,0);
  % Axes
  \draw[->,thick] (-3.5,0) -- (3.5,0) node[right] {$x$};
  \draw[->,thick] (0,-3.5) -- (0,0.5) node[above] {$y$};
  % Propagation along x
  \draw[arr,very thick] (-1.5,-0.3) -- (1.5,-0.3) node[pos=0.5,above,lbl] {propagation along $x$};
  % Decay with depth
  \draw[arr,very thick] (1.0,-0.6) -- (1.0,-2.6) node[pos=0.5,right,lbl] {decay with depth};
  % Elliptical particle motion at surface
  \draw[dashed] (-1.2,-0.4) to[out=60,in=120] (0.6,-0.2) to[out=-60,in=120] (2.4,-0.6);
  \node[node] at (0.4,0.6) {Rayleigh particle\\[-1pt] motion (elliptical)};
  \draw[arr] (0.4,0.6) -- (0.9,-0.1);
  % Region labels
  \node[above] at (0,0.3) {half-space $y>0$};
  \node[below] at (0,-3.2) {decay $e^{-m y}$};
\end{tikzpicture}

\caption{Geometry after \citep{Rayleigh1885}: thermoelastic Rayleigh wave on a half-space -- surface-trapped particle motion decaying exponentially with depth, coupled dilatational--thermal and shear potentials.}\label{fig:rayleigh-geo}
\end{figure}

\subsection{Results}
\Cref{fig:rayleigh} shows the speed change and damping of the thermoelastic Rayleigh wave for $\varepsilon=0.05$, $\nu=0.3$. The speed increases by about $2\times10^{-3}$ at low wavenumber relative to the purely elastic Rayleigh speed (the adiabatic limit) and tends to the elastic value at large $k$ (the isothermal limit), in the same way as the bulk wave (\Cref{fig:disp}a), but the effect is about an order of magnitude smaller than for the compressional wave because only part of the surface motion is dilatational. The damping peaks at $k\,c_1/\omega_*\approx0.4$ with a quality-factor reciprocal of about $2\times10^{-3}$. The thermal boundary condition matters at low wavenumber: an isothermal surface shifts the damping peak to smaller $k$ and lowers the speed change, since the surface temperature fluctuation is suppressed. The Lord--Shulman law ($\tau_q=0.25/\omega_*$, insulated surface) changes the results little at low wavenumber but increasingly at large $k$: relative to the classical insulated surface the damping differs by less than $1\,\%$ for $k\,c_1/\omega_*\le0.1$, by $5\,\%$ at the damping peak, by $10\,\%$ at $k\,c_1/\omega_*=1$ and saturates at about $12\,\%$ for $k\,c_1/\omega_*\ge10$; the speed offset is $-1.2\times10^{-4}$ in relative terms at large $k$. This is the pattern expected from the bulk wave: the heat law matters once the product of the frequency and the lag is of order one.

\begin{figure}[t]\centering
\includegraphics[width=\textwidth]{figs/fig_rayleigh.pdf}
\caption{Thermoelastic Rayleigh wave in a half-space ($\nu=0.3$, $\varepsilon=0.05$). (a) Speed relative to the purely elastic Rayleigh speed. (b) Temporal damping $Q^{-1}=2|\Imag\omega|/\Real\omega$ at real wavenumber $k$ in units of $\omega_*/c_1$. Curves: classical theory with insulated and with isothermal surface, and Lord--Shulman with $\tau_q=0.25/\omega_*$ and insulated surface.}\label{fig:rayleigh}
\end{figure}

\subsection{Anisotropic, piezoelectric and nonlocal surface waves}
For anisotropic media the Stroh formalism \citep{Stroh1958,Stroh1962} gives the displacement--traction state vector as a first-order system in $y$, and the secular equation becomes the vanishing of a $3\times3$ (or larger, when piezoelectric and thermal fields are added) determinant built from the eigenvectors with decaying exponents; this is the route taken for piezothermoelastic layers and half-spaces \citep{Biswas2020}. For piezoelectric materials, new surface-wave types arise: the Bleustein--Gulyaev shear-horizontal wave exists for transversely isotropic piezoelectric substrates even though a pure SH wave cannot exist at an ordinary elastic free surface; Rayleigh waves in nonlocal piezoelastic half-spaces and piezo-thermoelectric semiconductors have been treated with Eringen-type nonlocality and with fractional or Klein--Gordon-type extensions \citep{Tung2019,Gupta2024,Aldandani2026,Lazar2022}. Energy aspects of surface thermoelastic waves are discussed in \citep{IvanovSavova1991}. In each of these extensions the structure of the problem is identical to the determinant above, with additional exponents and boundary conditions; the physical effects are the shift of the Rayleigh root by electromechanical stiffening (by $K^2$, see Section~\ref{sec:multi}), by nonlocal dispersion (Section~\ref{sec:size}), and by thermal damping as shown here.

\subsection{Practical remarks on secular equations}
Three technical points recur. (i) The decaying branch of each square root must be chosen with $\Real m>0$, and in the complex plane this introduces branch cuts across which the determinant is not analytic; root finding that ignores this can jump to spurious roots. (ii) The determinant has trivial zeros ($s=0$ for instance) that must be excluded when tracking roots; we found that a naive secant search from a poor initial guess converged to such a trivial root. (iii) Continuation in $k$ from a reliable starting point (here the isothermal Rayleigh root) is more robust than independent root finding at each $k$.

\begin{summarybox}[Summary]
The thermoelastic Rayleigh problem is a $3\times3$ determinant \mathcal D(s,k) from two coupled dilatational--thermal modes and one shear mode, with insulated or isothermal surface. Speed shifts from adiabatic (low $k$) to elastic/isothermal (high $k$); damping peaks near $k c_1/\omega_*\approx0.4$. The thermal boundary condition shifts the peak; LS changes results by $\sim$10--12\% at large $k$.
\end{summarybox}
